\begin{abstract}
\noindent
Quantum reservoir computing (QRC) trains only a linear read-out on the measured
observables of a fixed, input-driven quantum system, sidestepping the trainability
problems of variational quantum circuits while inheriting the temporal-processing
strengths of classical reservoir computing. Renewable-energy and climate time
series --- solar irradiance, wind power, electric load, and coupled ocean--atmosphere
indices --- are a natural application domain: they are abundantly available, exhibit
memory and nonlinearity on many time scales, and support an economically important
forecasting industry against which any quantum claim can be honestly benchmarked.
This document is both a formal research review and an implementation handbook. It
develops classical reservoir computing and its guarantees (echo-state property,
fading memory, universality); formulates QRC in the language of quantum channels
and derives where nonlinearity, memory, and measurement enter; surveys the
principal reservoir architectures and hardware platforms; provides complete,
line-by-line-explained implementations in Cirq, Qiskit, and PennyLane, including
transpilation to native gate sets, realistic noise, finite-shot estimation, and
mid-circuit measurement with reset; catalogues and analyses public renewable-energy
and climate datasets with full exploratory data analysis of three of them; assembles
end-to-end forecasting pipelines that share one mathematical structure across
frameworks; and closes with design principles and ten concrete research directions.
All numerical results in this document were produced by the code shipped alongside
it, every experiment reports size-matched classical baselines and structureless
quantum controls, and no claim of quantum advantage is made where none was observed.
\end{abstract}

\vspace{2mm}
\subsection*{How to read this document}

Parts~I and~II build the theory in a linear order and are prerequisites for
everything else. Part~III (architectures) and Part~V (hardware) are reference
material: they can be read selectively. Part~IV is the software handbook and is
best read at a keyboard, alongside the \code{code/} directory that ships with
this document. Part~VI introduces the application domain and its data; Part~VII
assembles the pieces into complete pipelines; Part~VIII distils design judgement
and proposes research programmes. A reader who already knows classical reservoir
computing can start at Part~II; a reader who wants to forecast something by
tonight can start at Part~VII and backtrack as needed. Parts~IX--XI, added in
the July 2026 revision, are forward-looking: Part~IX designs the staged
repository that consolidates this document's code with the qubit-reuse
compiler, Part~X gives the theory and pre-registered implementation plan for
the recurrence-free reservoir with virtual-node readout, and Part~XI wraps
the forecast core in a full product pipeline with a survey of classical
ML/DL and agentic-AI post-processing.

\subsection*{Epistemic conventions}

Three kinds of statement appear in this document and are kept explicitly
distinct. \emph{Theory} means a mathematical statement with stated assumptions,
attributed to the literature or derived in the text. \emph{Simulation} means a
numerical result produced by the exact density-matrix or statevector code in
\code{code/}, on either synthetic surrogate data (always labelled
\emph{surrogate}) or real recorded data (always labelled \emph{real}).
\emph{Hardware} means a result reported in the literature from a physical
quantum device; this document runs no hardware jobs itself, but Part~IV's Qiskit
pipeline is written exactly as a hardware submission would be, including
transpilation, native basis gates, coupling constraints, realistic noise, and
finite shots. Throughout, reported errors are normalised mean-square errors
(NMSE), for which predicting the training-set mean scores $1.0$; any method
worth attention must sit well below that line, and every figure that shows an
NMSE also shows it.

\clearpage
\tableofcontents

\clearpage
\section*{Notation}
\addcontentsline{toc}{section}{Notation}

\begin{center}
\small
\begin{tabular}{@{}ll@{}}
\toprule
Symbol & Meaning \\
\midrule
$u_t \in \R^{d_{\mathrm{in}}}$ & input at discrete time $t$ (scalar unless stated) \\
$y_t$, $\hat y_t$ & target and predicted output \\
$x_t \in \R^{N_x}$ & classical reservoir state / extracted feature vector \\
$W_{\mathrm{res}}, W_{\mathrm{in}}, b$ & fixed reservoir, input, and bias weights \\
$W_{\mathrm{out}}$ & trained linear read-out \\
$f(\cdot)$ & elementwise activation (typically $\tanh$) \\
$\alpha$ & leak rate of a leaky-integrator reservoir \\
$\rho_t$ & reservoir density matrix at time $t$ \\
$\Hil$, $d=2^n$ & Hilbert space and its dimension for $n$ qubits \\
$\map{E}_{\theta,u}$ & input-dependent CPTP map (quantum reservoir update) \\
$K_j$ & Kraus operators, $\sum_j K_j^\dagger K_j = \id$ \\
$\map{L}$ & Lindblad generator \\
$O_k$, $x_t^{(k)}=\tr(O_k\rho_t)$ & measured observable and feature \\
$X_i, Y_i, Z_i$ & Pauli operators on qubit $i$ \\
$H$, $J_{ij}$, $h_i$ & Hamiltonian, couplings, fields \\
$\Delta t$ & input injection interval (evolution time per step) \\
$V$ & number of virtual nodes (temporal multiplexing) \\
$S$ & number of measurement shots \\
$L$ (or $W$) & input-window length in restart protocols \\
$\lambda$ & ridge-regression penalty \\
$\mathrm{MC}$, $\mathrm{MC}_d$ & total and delay-$d$ linear memory capacity \\
$\NMSE$ & mean-square error normalised by target variance \\
$\ev{r}$ & mean level-spacing ratio (chaos diagnostic) \\
$T_1$, $T_2$ & energy relaxation and dephasing times \\
\bottomrule
\end{tabular}
\end{center}
